\section{Original rectifications and harmonic families}
\subsection{Notation and the verified original construction}
The field $k$ has characteristic zero. We normalize the Chebyshev polynomials by
\[
C_j(u)=2T_j(u/2),\qquad C_0=2,\quad C_1=u,\quad C_{j+1}=uC_j-C_{j-1}.
\]
We consider closed embeddings, that is, $k[f_1,f_2,f_3]=k[u]$.
A parameter polynomial $P$ satisfies $P\circ\gamma=u$; to be a coordinate,
it must additionally extend to a polynomial automorphism. These two
conditions are not equivalent.

The original preprint \cite{Repository} provides automorphisms
$\Phi_T,\Phi_H$ and their polynomial inverses over
$K=\Q[a]/(a^2-a+1)\simeq\Q(\sqrt{-3})$ for
\[
\begin{aligned}
\gamma_T(u)&=(u^3-3u,u^4-4u^2,u^5-10u),\\
\gamma_H(u)&=(C_3(u),C_5(u),C_7(u)).
\end{aligned}
\]
The identities $\Phi_T\gamma_T=\Phi_H\gamma_H=(u,0,0)$, both inverse
compositions, and all twelve printed components have been verified by
exact polynomial identities. The component degrees are respectively
$(5,15,24)$ and $(7,21,41)$; their minimality is not asserted.
The complete formulas are retained in the original preprint and in the
self-contained computational package. The change of coordinates
\[
M(x,y,z)=(x,y+2,z-5x)
\]
takes $\gamma_T$ to $(C_3,C_4,C_5)$.

\subsection{Constructive transfer through a modification}
Let $A=(x,H,W)$ and $A^{-1}=(x,\mathcal Y,\mathcal Z)$ be
mutually inverse automorphisms, with $m\in k[x]$, $\deg m>0$, $p\in k[x,y]$.
Consider $\beta=(x,p(x,y)+m(x)z,y)$. In the quotient ring
$S=k[x]/(m)$, set $c(y)=H(x,p(x,y),y)$.
If the linear coefficient of $c$ is a unit and all higher coefficients
are nilpotent, the substitution $y\mapsto c(y)$ defines an automorphism of $S[y]$.
Its compositional inverse $\psi$ is obtained by a finite lifting process
along the nilpotent ideal generated by the higher coefficients.

Choose a polynomial representative
\[
\varphi(x,T)\equiv W(x,p(x,\psi(T)),\psi(T))\pmod m.
\]
Then
\[
\mathcal A=\left(x,H\circ\beta,
\frac{W\circ\beta-\varphi(x,H\circ\beta)}m\right)
\]
and
\[
\mathcal A^{-1}(x,T,Q)=
\left(x,\mathcal Z(x,T,mQ+\varphi),
\frac{\mathcal Y(x,T,mQ+\varphi)-p(x,\mathcal Z(x,T,mQ+\varphi))}m\right)
\]
have polynomial components: both numerators are divisible by the full $m$,
including the multiplicities of its roots. After inverting $m$, both
compositions are the identity. Injectivity of localization carries these
identities back to the original ring, including the exceptional fibre $m=0$.
This explicitly implements classical lifting and relative-coordinate
mechanisms; the general transfer idea is not claimed as new \cite{EVM,Repository}.
The case $m\in k^*$ is simpler: $\beta$ itself is already an automorphism.

If $H\circ\gamma=u$ and $w=W\circ\gamma$, the relative pair yields
a complete rectifying map
\[
\Phi=(H,\ x-f_1(H),\ W-w(H)).
\]
A parameter polynomial alone, without a verified coordinate pair,
does not suffice for this conclusion.

\subsection{Polynomial changes of harmonic indices}
For coprime $a,b$ and signs $\varepsilon_a,\varepsilon_b\in\{\pm1\}$,
the congruences
\[
c'\equiv\varepsilon_a c\pmod{2a},\qquad
c'\equiv\varepsilon_b c\pmod{2b}
\]
imply
\[
C_{c'}-\varepsilon_a\varepsilon_b C_c\in\Q[C_a,C_b].
\]
Consequently, the corresponding space curves are related by
a polynomial shear of the third coordinate.

A short proof uses the nodes of the plane projection $(C_a,C_b)$.
Its coordinate ring consists exactly of the polynomials in $u$ that have
equal values at the two preimages of every node: inclusion is immediate,
and equality follows by comparing codimensions
$(a-1)(b-1)/2$ through the semigroup $\langle a,b\rangle$ and interpolation.
The stated congruences of indices give the required equalities of values
by the trigonometric Chebyshev formula. Eliminating leading terms produces
a shear polynomial with rational coefficients. Details and exact examples
are given in the original appendix; the nodal projection is treated in \cite{FF}.

This gives two complete families over $K$:
\[
H(3,4,c),\quad\gcd(c,12)=1;
\qquad H(3,5,c),\quad\gcd(c,15)=1.
\]
For the pair $(3,4)$, the representatives $c=1,5$ modulo $24$ suffice;
for $(3,5)$, they are $c=1,2,4,7$ modulo $30$. Representatives with index
$1$ are immediate, those with index $2$ are classically known, and $c=4$
reduces by a permutation to $(3,4,5)$. The remaining two are supplied
by the original rectifications.

Index $1$, the admissible models $H(2,b,c)$ and $H(a,b,a+b)$ were already
known in the literature \cite{FF}; they are not attributed to this work.
The topological classification of harmonic knots \cite{KP} is likewise
not a classification of complex embeddings up to automorphisms.

\subsection{Finite enumeration and its precise scope}
All $44142$ pairwise coprime triples
$2\le a<b<c\le100$ have been enumerated using every strictly decreasing
elementary reduction
\[
c=ra+sb\quad\longmapsto\quad |ra-sb|,
\qquad (x,y,z)\longmapsto(x,y,C_r(x)C_s(y)-z).
\]
Of these, $3209$ reach index $1$, another $237$ reach a verified base
model, and $40696$ remain unresolved \emph{by this algorithm}.
The last number is not a count of nonrectifiable curves.

For any embedding $\gamma=(f_1,\ldots,f_n)$, we subsequently use
\[
A=k[t,u,X_1,\ldots,X_n],\qquad
D_\mu=t^\mu\partial_u+\sum_i f_i'(u)\partial_{X_i},\qquad
V=\gamma(u)-t^\mu X,
\]
where $\mu\ge1$. This is an LND in $n+2$ variables. It must not be confused
with the Jacobian LND of a modification in three variables.
